\scanpage{78}
% archive foliation: 77
4) \uline{Application aux immersions, etc}

\deleted{\uline{Proposition}}. Supposons que $W$ soit
\deleted{\unsure{affine} \unsure{spécial} \unsure{de type} $\unsure{\mathrm{II}}$}
\inserted{$\mathrm I'$, ou Greenberg spécial}.
Soit $f:\mathfrak X\to\mathfrak Y$ un morphisme d'immersion
(\deleted{\unsure{ouv.}} fermée).
\inserted{\deleted{loc. de type fini}}
Supposons \deleted{\unsure{en de}} \deleted{\unsure{que}}
dans le cas $\mathrm I'$ que
\inserted{tte partie finie d'un $\mathfrak Y_s$ soit contenue dans un ouvert affine de $\mathfrak Y$.}
\deleted{$\underline W^a(\mathfrak Y)$ existe et est réunion des
$\underline W^a(\mathcal U)$, $\mathcal U$ affines de $\mathfrak Y$.}
\deleted{Alors $\underline W^a(\mathfrak X)$ existe,}
et $\underline W^a(f)$ \inserted{(existe et)} est une immersion fermée.

En effet, on est ramené au cas affine, où cela a été vu.

\uline{Corollaire}
\inserted{(sous les conditions précédentes sur $W$)}
Si $\mathfrak Y$ est quasi-projectif sur $\unsure{S^*|S'}$,
et si $f:\mathfrak X\to\mathfrak Y$ est une immersion,
alors $\underline W^a(f)$ est une immersion.

En effet, $f$ se factorise en
\diagram{\begin{tikzcd}
\mathfrak X \arrow[r,"f'"] & \mathfrak X' \arrow[r,"f''"] & \mathfrak Y
\end{tikzcd}}
\deleted{\unsure{imm.}} $f'$ une immers. fermée et $f''$ imm.
ouverte. $\mathfrak X'$ est quasi-projectif \unsure{sur}
